\begin{lemma}[Uniform convergence of the blow-ups on compact sets]
Let $\Omega\subseteq\mathrm{Plane}$ be open, let $u\colon\mathrm{Plane}\to\mathbb{R}$, let
$x_0\in\Omega$ and let $m\in\mathbb{N}$ with $m\ge1$.  Assume that $u$ is of class $C^{m}$ on
$\Omega$ and that
\[
  D^{j}u(x_0)[w,\ldots,w]=0\qquad\text{for every }j\text{ with }1\le j<m
  \text{ and every }w\in\mathrm{Plane}.
\]
For $r>0$ and $x\in\mathrm{Plane}$ put
\[
  u_r(x):=\frac{u(x_0+r\,x)-u(x_0)}{r^{m}},\qquad
  P_m(x):=\frac{1}{m!}\,D^{m}u(x_0)[x,\ldots,x],
\]
where $D^{j}u(x_0)$ is the $j$-th Fr\'echet derivative of $u$ at $x_0$ and the brackets denote
its evaluation on the indicated number of copies of the vector shown.  Then for every $R>0$ the
family $u_r$ converges to $P_m$ uniformly on the closed ball $\overline{B_R(0)}$ as $r$
decreases to $0$ through positive values; explicitly, for every $R>0$ and every
$\varepsilon>0$ there is $\delta>0$ such that
\[
  \lvert u_r(x)-P_m(x)\rvert<\varepsilon
  \qquad\text{for all }r\text{ with }0<r<\delta\text{ and all }x\text{ with }\lvert x\rvert\le R.
\]
\end{lemma}
\begin{proof}
Fix $R>0$ and $\varepsilon>0$.

\medskip\noindent\emph{Step 1: a ball inside $\Omega$, and continuity of $D^{m}u$ at $x_0$.}
Since $\Omega$ is open and $x_0\in\Omega$, there is $\rho>0$ with
$\overline{B_\rho(x_0)}\subseteq\Omega$.  Since $u$ is of class $C^{m}$ on the open set
$\Omega$, clause (5) of the background proposition on the recursion for the iterated Fr\'echet
derivative says that the map $y\mapsto D^{m}u(y)$, with values in the normed space of
continuous $m$-linear forms on $\mathrm{Plane}$, is continuous on $\Omega$.  Put
\[
  \eta:=\frac{\varepsilon\,m!}{2R^{m}}>0 .
\]
By continuity of $y\mapsto D^{m}u(y)$ at $x_0\in\Omega$ there is $\delta_0$ with
$0<\delta_0\le\rho$ such that
\begin{equation}\label{eq:buc-cont}
  \bigl\lVert D^{m}u(y)-D^{m}u(x_0)\bigr\rVert\le\eta
  \qquad\text{whenever }\lvert y-x_0\rvert\le\delta_0 ,
\end{equation}
where $\lVert\cdot\rVert$ is the operator norm; note that any such $y$ lies in
$\overline{B_\rho(x_0)}\subseteq\Omega$, so $D^{m}u(y)$ is defined.  Set
\[
  \delta:=\frac{\delta_0}{R}>0 .
\]

\medskip\noindent\emph{Step 2: the one-variable restriction.}
Fix $r$ with $0<r<\delta$ and fix $x$ with $\lvert x\rvert\le R$.  Put
\[
  \gamma(s):=x_0+s\,x,\qquad \phi(s):=u(\gamma(s))\qquad(s\in\mathbb{R}),
\]
and let $J:=(-\delta,\delta)$, an open interval.  For $s\in J$,
\[
  \lvert\gamma(s)-x_0\rvert=\lvert s\rvert\,\lvert x\rvert\le\lvert s\rvert\,R<\delta R=\delta_0\le\rho ,
\]
so $\gamma(s)\in\overline{B_\rho(x_0)}\subseteq\Omega$.  Thus $\gamma(J)\subseteq\Omega$, and
$J$ contains the closed interval $[0,r]$ because $0\le r<\delta$.  The path $\gamma$ is affine
with $\gamma'(s)=x$, and $u$ is of class $C^{m}$ on the open set $\Omega$, so the chain-rule
background proposition gives that $\phi=u\circ\gamma$ is of class $C^{m}$ on $J$.

\medskip\noindent\emph{Step 3: the derivatives of $\phi$ at $0$.}
Let $0\le j\le m$.  Since $u$ is of class $C^{m}$ on $\Omega$ it is of class $C^{j}$ on
$\Omega$, and $\gamma(0)=x_0\in\Omega$.  Hence \noderef{IteratedDerivAlongRay}, applied with
the open set $U=\Omega$, the function $u$, the order $j$, the base point $x_0$, the direction
$x$ and the parameter $t=0$, gives
\begin{equation}\label{eq:buc-at0}
  \phi^{(j)}(0)=D^{j}u(x_0)[x,\ldots,x]\qquad(0\le j\le m).
\end{equation}
For $j=0$ this reads $\phi^{(0)}(0)=u(x_0)$.  For $1\le j\le m-1$ the hypothesis of the lemma,
taken with $w=x$, makes the right-hand side of \eqref{eq:buc-at0} vanish, so
\begin{equation}\label{eq:buc-vanish}
  \phi^{(j)}(0)=0\qquad\text{for }1\le j\le m-1 .
\end{equation}
(For $m=1$ the range $1\le j\le m-1$ is empty and \eqref{eq:buc-vanish} is vacuous.)

\medskip\noindent\emph{Step 4: Taylor's theorem of order $m-1$ along the ray.}
Apply the Taylor--Lagrange background proposition with $n:=m-1$ --- legitimate because $m\ge1$,
so $n\in\mathbb{N}$ and $n+1=m$ --- with $a:=0$ and $b:=r$, which are distinct since $r>0$, and
with the function $\phi$, which by Step 2 is of class $C^{n+1}=C^{m}$ on the open interval
$J\supseteq[0,r]$.  It produces $\xi$ strictly between $0$ and $r$, so $\xi\in(0,r)\subseteq J$,
with
\[
  \phi(r)=\sum_{j=0}^{m-1}\frac{\phi^{(j)}(0)}{j!}\,r^{j}
          +\frac{\phi^{(m)}(\xi)}{m!}\,r^{m} .
\]
By \eqref{eq:buc-vanish} all the terms with $1\le j\le m-1$ vanish, and the $j=0$ term is
$\phi(0)=u(x_0)$.  Hence
\begin{equation}\label{eq:buc-taylor}
  u(x_0+r\,x)=\phi(r)=u(x_0)+\frac{r^{m}}{m!}\,\phi^{(m)}(\xi).
\end{equation}
Since $\gamma(\xi)\in\Omega$ by Step 2, \noderef{IteratedDerivAlongRay} applies again, now with
the order $m$ and the parameter $t=\xi$, and gives
\begin{equation}\label{eq:buc-atxi}
  \phi^{(m)}(\xi)=D^{m}u(x_0+\xi\,x)[x,\ldots,x].
\end{equation}

\medskip\noindent\emph{Step 5: the estimate.}
Dividing \eqref{eq:buc-taylor} by $r^{m}>0$ and inserting \eqref{eq:buc-atxi},
\[
  u_r(x)=\frac{u(x_0+r\,x)-u(x_0)}{r^{m}}
        =\frac{1}{m!}\,D^{m}u(x_0+\xi\,x)[x,\ldots,x].
\]
Subtracting $P_m(x)=\frac1{m!}D^{m}u(x_0)[x,\ldots,x]$ and using that evaluation at the fixed
family $(x,\ldots,x)$ is linear in the multilinear form,
\[
  u_r(x)-P_m(x)=\frac{1}{m!}\,\bigl(D^{m}u(x_0+\xi\,x)-D^{m}u(x_0)\bigr)[x,\ldots,x].
\]
The operator-norm background proposition therefore gives
\[
  \lvert u_r(x)-P_m(x)\rvert
  \le\frac{1}{m!}\,\bigl\lVert D^{m}u(x_0+\xi\,x)-D^{m}u(x_0)\bigr\rVert\;\lvert x\rvert^{m}.
\]
Now $\lvert(x_0+\xi\,x)-x_0\rvert=\xi\,\lvert x\rvert\le r\,R<\delta R=\delta_0$, using
$0<\xi<r$ and $\lvert x\rvert\le R$, so \eqref{eq:buc-cont} applies at $y=x_0+\xi\,x$ and bounds
the operator norm by $\eta$.  Together with $\lvert x\rvert^{m}\le R^{m}$ (monotonicity of
$t\mapsto t^{m}$ on $[0,\infty)$) this yields
\[
  \lvert u_r(x)-P_m(x)\rvert
  \le\frac{\eta\,R^{m}}{m!}
  =\frac{1}{m!}\cdot\frac{\varepsilon\,m!}{2R^{m}}\cdot R^{m}
  =\frac{\varepsilon}{2}<\varepsilon .
\]

\medskip\noindent Since $r\in(0,\delta)$ and $x$ with $\lvert x\rvert\le R$ were arbitrary, and
$\delta$ was produced from $R$ and $\varepsilon$ alone, this is exactly the asserted uniform
convergence on $\overline{B_R(0)}$ as $r$ decreases to $0$ through positive values.
\end{proof}
